\subsection{Mobility and transport impacts}

    Bike-sharing has measurably altered urban mobility patterns by adding a flexible, short-trip transport option that often complements existing modes. In general, \acp{bss} trips have been observed to be predominantly utilitarian, where commuting and errands are the most common, but they also support leisure and tourism mobility~\cite{Ricci2015}. 

    A widely mentioned benefit is improved first- and last-mile connectivity to public transportation~\cite{Shaheen2010, Qiu2018}. By strategically siting bike stations near public transport hubs, \ac{bss} can extend the reach of subways and buses, and therefore facilitate multimodal trips~\cite{Cheng2022}. In Beijing, for example, bike-sharing access to bus and metro stations saved commuters an average of eight minutes per day, translating into substantial citywide time savings and even modest productivity gains~\cite{Qiu2018}. Recent analyses further highlight efficiency gains from e-bikes, which achieve travel speeds roughly 9\% higher than mechanical bikes, equivalent to roughly 70 seconds saved on a 1.5--2 km trip~\cite{Waldner2025}, further highlighting the contribution of contemporary \acp{bss} to urban time efficiency.  

    Beyond connectivity, \acp{bss} capacity for modal substitution is central, as it is one of the most debated aspects of their impact. Understanding which modes shared bikes replace is essential, as reducing car dependency is a core motivation for their introduction~\cite{Fishman2016}. Early surveys of the 2010s across cities including Dublin, London, Lyon, Barcelona, Montreal, Washington D.C., Minneapolis, Melbourne and Brisbane consistently showed that only 5--20\% of bike-share trips would otherwise have been made by car, with the majority substituting for walking or public transport~\cite{Murphy2015, Fishman2013, Fishman2014, Bachand-Marleau2012}. Ricci~\cite{Ricci2015} therefore concluded that there was “no evidence that bike sharing significantly reduces traffic congestion, carbon emissions and pollution”. At the same time, she also cautioned that mode substitution alone does not capture the full range of impacts, since surveys rarely measure the characteristics of replaced car trips (e.g. frequency, duration, route), or the extent of induced demand through new trips not previously made. On this basis, both Ricci and Fishman called for more detailed research on this issue on 2015~\cite{Ricci2015, Fishman2016}.  

    Subsequent research broadly supports these findings. Despite the technological advances of \ac{e-b} and dockless systems, the consensus remains the same: while bike-sharing does replace some car usage, the majority of trips continue to come from walking and public transport. In Beijing a survey of dockless system users found that only 14.8\% of trips replaced car travel~\cite{Chen2022b}; in Delft, the introduction of both docked (\textit{OV-fiets}) and dockless (\textit{Mobike}) systems showed similar patterns, with approximately 30\% of trips replacing car use~\cite{Ma2020}; and a multi-city European survey in 2021--2022 placed the figure at about 21\%~\cite{Bosehans2024}. Although these figures are higher than early estimates, still indicate that most bike-share usage comes at the expense of active modes or transit rather than driving. Accordingly, some authors argue that more restrictive measures, such as car-free zones or increased parking charges, are necessary to strengthen the car-reducing effects of \acp{bss}~\cite{Bosehans2024}.    

    Context, however, matters. Research highlights that \ac{bss} user demographics and trip purposes shape modal substitution outcomes. Dockless systems in China and Europe often attract younger riders and those without cars or driving licenses, meaning these users were unlikely to be replacing car trips in the first place~\cite{Ma2020, Bosehans2024}. By contrast, in cities or niches where bike-share directly competes with car convenience, higher car substitution rates are possible. Moreover, cargo bike-sharing services---which allow users to haul goods or children---show particularly strong potential to displace car use: one survey of cargo bike-share users found that 46\% would have driven a car if the cargo bike had not been available~\cite{Becker2018}. This suggests that tailoring bike-share offerings (e.g. e-bikes for longer commutes, cargo bikes for errands) can increase the likelihood of replacing car trips.

    Although direct car substitution rates remain modest, even small shifts can yield measurable system-level benefits. In Washington D.C., \textit{Capital Bike-share} reduced neighborhood traffic congestion by about 4\%, with the greatest reductions in already congested areas~\cite{Hamilton2018}. A study of 96 U.S. cities similarly found that bike-share deployment eased congestion in large, transit-rich cities, particularly at rush hours, though effects were weaker or mixed in smaller and high-car-ownership contexts~\cite{Wang2017}. In China’s megacities such as Beijing, Shanghai, and Wuhan, the initial rollout of dockless bike-share significantly reduced congestion in the short term, though subsequent expansions had diminishing returns; notably, metro ridership rose alongside bike-share use, suggesting that many users shifted from driving to bike+metro combinations~\cite{Xu2024}. Parking impacts follow a similar logic: in Pittsburgh, the opening of a bike-share station was associated with a 2\% decline in parking meter use~\cite{Pelechrinis2016}, while in Boston the installation of new stations led to a 2.2\% reduction in household vehicle ownership and a 3.3\% drop in vehicle-miles traveled, with the strongest impacts (around 10\% lower car dependence) observed near transit stations~\cite{Basu2021}. Together, these findings suggest that \acp{bss}, as part of a multimodal system, can support “car-lite” lifestyles and free up curb space over the long term.  

    Beyond substitution, \acp{bss} also induce new trips that would not otherwise occur. A European study found that about 11\% of shared mobility trips were entirely new journeys induced by the service~\cite{Bosehans2024}. Such trips do not reduce any particular mode, but they indicate that bike-sharing provides additional mobility and access, expanding travel opportunities for its users. Moreover, even when not replacing car journeys outright, bike-share can work in tandem with public transport to reduce car dependence. Many users combine bike-share with public transport for first- and last-mile travel, and surveys report that participation in bike-share often leads to a reduction in car use frequency. In Washington D.C., for example, 55\% of respondents to Washington D.C.'s \textit{CaBi} 2016 annual survey reported driving less often since joining the system~\cite{NACTO2020}, underscoring how bike-share can gradually chip away at car use through a mix of direct substitution and enabling multimodal lifestyles.

    Finally, another key mobility impact of \acp{bss} is their contribution to the transport network resilience. During four London metro strikes, bike-share trips spiked sharply, with many people who would have used the metro switching to shared bikes, often using longer or unfamiliar routes when their regular stations were inaccessible~\cite{Yang2022}. Similarly, during a major flood in Porto Alegre, Brazil, the local bike-share system continued to function and enabled many daily activities to proceed despite widespread transport disruption~\cite{Araujo2025}. Evidence from Washington D.C. further shows that bike-share usage rises significantly during planned Metro shutdowns, demonstrating its role as an alternative mode in complete transit outages~\cite{Cheng2022}. Last, the COVID-19 pandemic further reinforced this role of \acp{bss} as a transport mode that brings resilience to the transport network. While public transport usage remained far below pre-pandemic levels, bike-sharing volumes nearly returned to normal once restrictions eased, with some even increasing its trips numbers~\cite{Teixeira2024}. For instance, in Barcelona's \textit{Bicing}, trip levels exceeded pre-pandemic counts within just a couple of months of reopening after the quarantine~\cite{Grau-Escolano2024}.